% Generated draft: requires mathematical and source review.
\documentclass[11pt]{article}
\usepackage[margin=25mm]{geometry}
\usepackage{fontspec}
\setmainfont{Latin Modern Roman}
\usepackage{amsmath,amssymb,mathtools,unicode-math}
\setmathfont{Latin Modern Math}
\usepackage{enumitem,needspace,xurl,hyperref}
\setlength{\parindent}{0pt}
\setlength{\parskip}{5pt}
\setlength{\emergencystretch}{3em}
\newcommand{\N}{\mathbb{N}}
\newcommand{\Z}{\mathbb{Z}}
\newcommand{\Q}{\mathbb{Q}}
\newcommand{\R}{\mathbb{R}}
\newcommand{\C}{\mathbb{C}}
\newcommand{\F}{\mathbb{F}}
\newcommand{\A}{\mathbb{A}}
\newcommand{\PP}{\mathbb{P}}
\newcommand{\E}{\mathbb{E}}
\newcommand{\eps}{\varepsilon}
\newcommand{\dd}{\,\mathrm{d}}
\begin{document}
\section*{UnsolvedMath ID: 30005643}
\textbf{Source problem code:} catalogue entry
\subsection{Definitions and mathematical statement}
Let \(q=p^r\) be a prime power and let \(K=\F_q(t)\). Let \(G\leq S_m\) be a transitive permutation group. The supplied statement asks for a lower bound of Malle type for a counting function \(N_{G,K}(X)\) of degree-\(m\) extensions with discriminant norm at most \(X\).

A conventional interpretation of the group condition is that \(L/K\) is finite separable and that the Galois closure \(\widetilde L/K\), acting on the \(m\) embeddings of \(L\) into a separable closure of \(K\), has permutation group conjugate to \(G\) in \(S_m\). This interpretation requires confirmation: the source says that the degree-\(m\) extension has Galois group \(G\), without explicitly mentioning its Galois closure. Requiring \(L/K\) itself to be Galois would give a different counting problem.

Under the Galois-closure interpretation, and once the discriminant and constant-field conventions have been fixed, the counting function would be
\[
N_{G,K}(X)=\#\left\{[L/K]_K:\ [L:K]=m,\ L/K\text{ finite separable},\ \operatorname{Gal}(\widetilde L/K)\cong G\text{ in the specified permutation action},\ \mathrm N\mathfrak d_{L/K}\leq X\right\},
\]
where \([L/K]_K\) denotes a \(K\)-isomorphism class. If \(\mathfrak d_{L/K}\) means the full discriminant divisor on \(\PP^1_{\F_q}\), then its norm is \(q^{\deg\mathfrak d_{L/K}}\). Using instead a discriminant ideal over \(\F_q[t]\) excludes the place at infinity; the supplied record does not choose between these conventions.

For comparison, for nontrivial \(G\) the usual permutation-index exponent is defined by
\[
\operatorname{ind}(g)=m-\#\{\text{orbits of }\langle g\rangle\text{ on }\{1,\ldots,m\}\},
\qquad
 a(G)=\left(\min_{1\ne g\in G}\operatorname{ind}(g)\right)^{-1}.
\]
A precise power-exponent lower-bound formulation would be
\[
\forall\varepsilon>0\ \exists c=c(G,q,\varepsilon)>0\ \exists X_0=X_0(G,q,\varepsilon)>0\ \forall X\geq X_0:\quad
N_{G,K}(X)\geq cX^{a(G)-\varepsilon}.
\]
This formula is a possible interpretation to be checked, not an identification established by the supplied extract. A lower bound retaining the conjectured logarithmic factor is a stronger possible target. The record supplies neither the exact lower-bound formula nor the logarithmic exponent.

\textbf{Unresolved statement qualifications.} The requested target is the conjectural lower bound in Malle's asymptotic for \(N_{G,\F_q(t)}(X)\). A definitive quantified statement cannot be recovered from the supplied material without specifying the strength of that bound, the admissible \(q\), any restrictions relating \(p\) to \(G\), whether the extensions must have constant field exactly \(\F_q\), and the discriminant convention. The trivial transitive group also requires separate treatment because the minimum defining \(a(G)\) is empty. No resolution of these issues is asserted here.
\subsection{Short English statement}
The supplied problem asks for the conjectured lower bound on the number of degree-\(m\) extensions of \(\F_q(t)\) with prescribed transitive Galois-closure group and bounded discriminant. The exact bound and several counting hypotheses are missing. The displayed power-exponent formulation is a possible interpretation requiring review. The supplied label is \emph{partially solved}; it is not an independently verified status.
\subsection{Sources}
\begin{itemize}
\item Supplied local UnsolvedMath registry record, permanent ID \(30005643\), title \emph{Malle Lower Bounds over Rational Function Fields}, source problem code \texttt{catalogue entry}. Project registry references: \texttt{data/unsolvedmath/problems.json} and \texttt{catalog/problems.json}. This draft uses the supplied record; no repository or external-source inspection was performed.
\item Supplied report citation: \emph{MFO--RIMS Tandem Workshop: Arithmetic Homotopy and Galois Theory (2024)}, Oberwolfach Reports, pp.~2377--2488, DOI \texttt{10.4171/owr/2023/42}. \url{https://doi.org/10.4171/owr/2023/42}. The record separately gives proposed year 2023. No individual proposer is supplied.
\item Supplied literature reference: Ellenberg--Tran--Westerland, \emph{Malle over function fields}. \url{https://arxiv.org/abs/1701.04541}. This is the label supplied by the registry; its bibliographic details and applicable hypotheses have not been independently checked.
\item Supplied literature reference: \emph{OWR report}. \url{https://doi.org/10.4171/owr/2023/42}.
\item Literature assessment dated exactly \texttt{2026-08-22}: the supplied curation labels the problem partially solved and states that the abelian case is a theorem, substantial nonabelian upper bounds are known over function fields, and the required lower bound remains missing in general. This is supplied dated curation, not independently verified mathematical or literature assessment; the scope and hypotheses of these claims require checking.
\item The registry labels the statement \texttt{corrected\_verified} and reports that it was checked against the official report. That verification is a supplied provenance claim. The supplied original extract instead discusses cohomological methods for arithmetic statistics, mentions Ellenberg and Landesman with label [EL23], Bhargava--Kane--Lenstra--Poonen--Rains heuristics with label [BKL+15], [EVW16], and Das with label [Das21], and ends mid-sentence. Complete bibliographic details for those labels are not supplied, and the extract does not state the requested Malle lower bound.
\end{itemize}
\end{document}
